\chapter{1977年辽宁卷}

\section{解答题}

\begin{problem}
计算 \( \sqrt{a^2-2ab+b^2} \) （\( a<b \)）.
\end{problem}

\begin{solution}
因为 \(a<b\)，所以 \(a-b<0\)，从而 \(\lvert a-b\rvert=b-a\). 因此原式
\(=\sqrt{a^2-2ab+b^2}=\sqrt{(a-b)^2}=\lvert a-b\rvert=b-a\).
\end{solution}

\begin{problem}
把 \( x^5y-x^3y+2x^2y-xy \) 分解因式.
\end{problem}

\begin{solution}
先提取公因式 \(xy\)，得原式 \(=xy\left(x^4-x^2+2x-1\right)\). 又因为
\(x^2-2x+1=(x-1)^2\)，所以 \(x^4-x^2+2x-1=x^4-(x-1)^2\). 这是平方差，故
\(x^4-(x-1)^2=\left(x^2+x-1\right)\left(x^2-x+1\right)\). 因此原式
\(=xy\left(x^2+x-1\right)\left(x^2-x+1\right)\).
\end{solution}

\begin{problem}
计算：\( \frac{0.027^{-\frac{1}{3}}+\left( -\frac{1}{6} \right)^{-2}-3^{-1}}{\left( -1000 \right)^0-\log_{10}\sqrt[3]{10}}\times 5^{\log_{5}\frac{2}{3}} \).
\end{problem}

\begin{solution}
因为 \(0.027=\left(\frac{3}{10}\right)^3\)，所以 \(0.027^{-\frac13}=\frac{10}{3}\)；又
\(\left(-\frac16\right)^{-2}=36\)，\(3^{-1}=\frac13\)，\((-1000)^0=1\)，且
\(\log_{10}\sqrt[3]{10}=\log_{10}10^{1/3}=\frac13\). 另外，
\(5^{\log_5\frac23}=\frac23\). 因而原式
\(=\frac{\frac{10}{3}+36-\frac13}{1-\frac13}\times\frac23
=\frac{39}{\frac23}\times\frac23=39\).
\end{solution}

\begin{problem}
如图，\(M\) 表示地平面，\(PC\) 表示直立地面的电柱，\( AC=6\mathrm{~m} \)，\(BC\) 是 \(AC\) 的 \( \frac{2}{3} \)，求拉线 \(AB\) 的长.

\begin{center}
\bitmapfigure[float=inline,max width=4.8cm]{img/1977/liaoning/q4_fig1.png}
\end{center}
\end{problem}

\begin{solution}
由 \(PC\perp\) 平面 \(M\)，且 \(BC\) 在平面 \(M\) 内，得 \(PC\perp BC\). 图中 \(A\) 在电柱 \(PC\) 上，因此 \(AC\perp BC\)，\(\triangle ABC\) 是以 \(C\) 为直角的直角三角形. 又
\(AC=6\mathrm{~m}\)，\(BC=\frac23AC=4\mathrm{~m}\). 由勾股定理，
\(AB^2=AC^2+BC^2=6^2+4^2=52\)，且 \(AB>0\)，所以
\(AB=\sqrt{52}=2\sqrt{13}\mathrm{~m}\). 即拉线 \(AB\) 的长为 \(2\sqrt{13}\mathrm{~m}\).
\end{solution}

\begin{problem}
求出 \( A(-3,3) \)，\( B(0,-1) \)，\( C(1,0) \) 三点中每两点间的距离，并指出以这三点为顶点的三角形 \(ABC\) 是什么样的三角形.

\bitmapfigure[max width=4.4cm]{img/1977/liaoning/q5_fig1.png}
\end{problem}

\begin{solution}
由两点间距离公式，
\(AB=\sqrt{[0-(-3)]^2+(-1-3)^2}=\sqrt{9+16}=5\)，
\(AC=\sqrt{[1-(-3)]^2+(0-3)^2}=\sqrt{16+9}=5\)，
\(BC=\sqrt{(1-0)^2+[0-(-1)]^2}=\sqrt2\). 因为 \(AB=AC\)，所以 \(\triangle ABC\) 为等腰三角形，且底边为 \(BC\).

\end{solution}

\begin{problem}
判别二次函数 \( y=x^2-4x+3 \) 有最大值还是有最小值？并把它求出来；写出这个函数图象的顶点坐标，用描点法画出这个函数的图象（画出草图即可）.
\end{problem}

\begin{solution}
【法一】二次函数的二次项系数为 \(a=1>0\)，所以抛物线开口向上，函数有最小值而无最大值. 对
\(y=x^2-4x+3\) 配方，得 \(y=(x-2)^2-1\). 因为 \((x-2)^2\geqslant0\)，所以当且仅当 \(x=2\) 时，
\(y\) 取得最小值 \(-1\)，没有最大值. 顶点坐标为 \((2,-1)\)，对称轴为 \(x=2\). 取若干对称点：
\(x=0\) 时 \(y=3\)，\(x=1\) 时 \(y=0\)，\(x=2\) 时 \(y=-1\)，\(x=3\) 时 \(y=0\)，\(x=4\) 时 \(y=3\). 因而在坐标系中描出
\((0,3)\)、\((1,0)\)、\((2,-1)\)、\((3,0)\)、\((4,3)\) 五点，并以平滑曲线连接，所得图象是开口向上、顶点为 \((2,-1)\)、对称轴为 \(x=2\) 的抛物线.

\begin{center}
\begin{tabular}{|c|c|c|c|c|}
\hline
\(x\) & \(\cdots\) & 0 & 1 & 2 \\
\hline
\(y\) & \(\cdots\) & 3 & 0 & \(-1\) \\
\hline
\end{tabular}
\end{center}

\solutionfigure{\bitmapfigure[max width=5.4cm]{img/1977/liaoning/q6_solution_fig1.png}}

【法二】由 \(a=1>0\)，函数有最小值. 配方得
\(y=x^2-4x+3=(x-2)^2-1\). 因为平方项不小于零，所以当且仅当 \(x=2\) 时，\(y\) 最小，最小值为 \(-1\)，且无最大值；顶点坐标为 \((2,-1)\). 由对称轴 \(x=2\)，只需在顶点两侧取对称点即可作出同一条抛物线. 例如取 \(x=1\)，\(2\)，\(3\)，分别得到 \(y=0\)，\(-1\)，\(0\)，再取 \(x=0\)，\(4\)，分别得到 \(y=3\)，\(3\). 描出这些点并用平滑曲线连接，即得到开口向上、顶点为 \((2,-1)\) 的抛物线.

\end{solution}

\begin{problem}
制造某种产品，计划经过两年使成本降低 \( 36\% \)，问平均每年应降低百分之几？
\end{problem}

\begin{solution}
设原来成本为 \(1\)，平均每年应降低的百分数为 \( x \)，根据题意，得
\( \left( 1-x \right)^2=1-36\% \).
\( \left( 1-x \right)^2=0.64 \).
由于每年降低的百分数满足 \(0\leqslant x<1\)，故 \(1-x>0\)，应取正平方根：
\(1-x=\sqrt{0.64}=0.8\)，所以 \(x=0.2=20\%\). 若列出平方方程的另一代数根，则 \(1-x=-0.8\) 给出 \(x=1.8=180\%\)，与降低百分数的实际范围不符，舍去.

答：平均每年应降低 \( 20\% \).

本题也可设原来成本为 \( a \)（\( a>0 \)），平均每年应降低的百分数为 \( x \)，根据题意，得 \( a\left( 1-x \right)^2=a\left( 1-36\% \right) \). 约去 \(a>0\) 后，仍得到
\(\left(1-x\right)^2=0.64\)，结合 \(0\leqslant x<1\) 可得 \(x=20\%\)，与前面的结果一致.
\end{solution}

\begin{problem}
在 \( \triangle ABC \) 中，已知 \( \cos A=\frac{5}{13} \)，\( \cos B=\frac{3}{5} \)，求 \( \cos C \) 的值.
\end{problem}

\begin{solution}
三角形内角满足 \(0^\circ<A<180^\circ\)、\(0^\circ<B<180^\circ\)，而余弦值为正，所以 \(A\)、\(B\) 均为锐角，正弦值取正根. 又因为 \(A+B+C=180^\circ\)，
\(C=180^\circ-(A+B)\)，故
\(\cos C=\cos[180^\circ-(A+B)]=-\cos(A+B)=\sin A\sin B-\cos A\cos B\).
由已知，\(\sin A=\sqrt{1-(5/13)^2}=12/13\)，\(\sin B=\sqrt{1-(3/5)^2}=4/5\). 因而
\(\cos C=\frac{12}{13}\cdot\frac45-\frac5{13}\cdot\frac35=\frac{48-15}{65}=\frac{33}{65}\).
\end{solution}

\begin{problem}
求证：两条对角线相等的梯形是等腰梯形（画出图形，写出：已知，求证，证明）.
\end{problem}

\begin{solution}
已知：四边形 \(ABCD\) 是梯形，\( DC \parallel AB \)，\(AC\) 和 \(BD\) 是对角线，且 \( AC=BD \).

求证：\( AD=BC \).

【法一】分别过 \(C\)、\(D\) 作 \(AB\) 的垂线，垂足为 \(E\)、\(F\). 因为 \(DC\parallel AB\)，且 \(CE\)、\(DF\) 都是两平行线 \(DC\)、\(AB\) 之间的距离，所以 \(CE=DF\). 又 \(AC=BD\)，且 \(\angle AEC=\angle BFD=90^\circ\)，故由斜边和一条直角边分别相等，得 \(\triangle ACE\cong\triangle BDF\). 从而 \(\angle CAE=\angle DBF\). 由于 \(A\)，\(E\)，\(B\) 共线、\(F\)，\(B\)，\(A\) 共线，得到 \(\angle CAB=\angle DBA\). 在 \(\triangle CAB\) 与 \(\triangle DBA\) 中，\(AC=BD\)、\(\angle CAB=\angle DBA\)、\(AB=BA\)，所以由边角边得 \(\triangle CAB\cong\triangle DBA\)，从而 \(AD=BC\).

\solutionfigure{\bitmapfigure[max width=6.0cm]{img/1977/liaoning/q9_solution_fig1.png}}

【法二】过 \(C\) 作 \(CE\parallel DB\)，交 \(AB\) 的延长线于 \(E\). 因为 \(DC\parallel AB\)，且 \(BE\) 在 \(AB\) 上，所以 \(DC\parallel BE\). 因此四边形 \(BECD\) 的两组对边分别平行，是平行四边形，从而 \(EC=BD\). 结合已知 \(AC=BD\)，得 \(AC=EC\)，所以 \(\triangle ACE\) 为等腰三角形，\(\angle CAE=\angle CEA\). 由于 \(A\)，\(E\)，\(B\) 共线，\(\angle CAE=\angle CAB\)；又因 \(CE\parallel DB\)，\(\angle CEA=\angle DBA\). 故 \(\angle CAB=\angle DBA\). 在 \(\triangle DAB\) 与 \(\triangle CBA\) 中，\(DB=AC\)、\(\angle DBA=\angle CAB\)、\(AB=BA\)，由边角边得两三角形全等，故 \(AD=BC\).

\solutionfigure{\bitmapfigure[max width=6.2cm]{img/1977/liaoning/q9_solution_fig2.png}}

【法三】设对角线 \(AC\)、\(BD\) 交于 \(O\). 由 \(DC\parallel AB\)，得 \(\angle 1=\angle 2\)、\(\angle 3=\angle 4\)，所以 \(\triangle AOB\sim\triangle COD\). 因此
\(\frac{AO}{OC}=\frac{BO}{OD}\). 两边分别加 \(1\)，得
\(\frac{AO+OC}{OC}=\frac{BO+OD}{OD}\)，即 \(\frac{AC}{OC}=\frac{BD}{OD}\). 由已知 \(AC=BD\)，可得 \(OC=OD\). 在等腰三角形 \(COD\) 中，等边对等角，所以 \(\angle 2=\angle 4\). 于是，在 \(\triangle ADC\) 和 \(\triangle BCD\) 中，\(AC=BD\)、\(DC=DC\)、\(\angle ACD=\angle BDC\)，由边角边得两三角形全等，故 \(AD=BC\).

\solutionfigure{\bitmapfigure[max width=6.2cm]{img/1977/liaoning/q9_solution_fig3.png}}
\end{solution}

\begin{problem}
\( \alpha \) 是多少度时，方程 \( 2x^2+2\sqrt{2}x+\tan \alpha=0 \) 有两个相等的实数根？（\( 0^\circ\leqslant \alpha<360^\circ \)）
\end{problem}

\begin{solution}
这是关于 \(x\) 的一元二次方程，二次项系数 \(2\neq0\). 它有两个相等的实数根，当且仅当判别式为零. 取 \(a=2\)、\(b=2\sqrt2\)、\(c=\tan\alpha\)，得
\(b^2-4ac=(2\sqrt2)^2-4\cdot2\tan\alpha=8-8\tan\alpha\). 因而
\(8-8\tan\alpha=0\)，即 \(\tan\alpha=1\). 在 \(0^\circ\leqslant\alpha<360^\circ\) 内，正切函数有定义且等于 \(1\) 的角为第一象限的 \(45^\circ\) 和第三象限的 \(225^\circ\). 所以所求为 \(\alpha=45^\circ\) 或 \(\alpha=225^\circ\).
\end{solution}

\begin{problem}
如图，我海军某舰艇在 \(A\) 处测得入侵敌舰在方位角为 \( 45^\circ \)，距离为 \(10\text{ 海里}\) 的 \(C\) 处，并正沿方位角为 \( 105^\circ \) 的方向以 \(9\text{ 海里/小时}\) 的速度逃跑，我舰艇立即以 \(21\text{ 海里/小时}\) 的速度沿直线前去追捕.

问我舰艇应沿多大的方位角和需用多少小时才能在某处 \(B\) 追上敌舰？

其中 \( \sin 21^\circ47'=\frac{3\sqrt{3}}{14} \)，\( \sin 38^\circ13'=\frac{5\sqrt{3}}{14} \).
这两个数据可根据解题需要取用.

\begin{center}
\bitmapfigure[float=inline,max width=5.8cm]{img/1977/liaoning/q11_fig1.png}
\end{center}
\end{problem}

\begin{solution}
设我舰艇在某处 \(B\) 追上敌舰需用 \(x\) 小时. 因为追捕需要正的时间，所以 \(x>0\). 在这段时间内，我舰艇和敌舰分别行驶 \(21x\)、\(9x\) 海里，故
\(AB=21x\text{ 海里}\)，\(BC=9x\text{ 海里}\)，而 \(AC=10\text{ 海里}\). 敌舰从 \(C\) 出发的方位角为 \(105^\circ\)，所以从 \(C\) 指向 \(A\) 的方向与从 \(C\) 指向 \(B\) 的方向夹角为
\(\angle C=45^\circ+(180^\circ-105^\circ)=120^\circ\).
由正弦定理，
\(\frac{21x}{\sin120^\circ}=\frac{9x}{\sin A}\). 由于 \(x>0\)，可约去 \(x\)，得
\(\sin A=\frac{9}{21}\sin120^\circ=\frac37\sin60^\circ=\frac{3\sqrt3}{14}\). 又因
\(A+B=60^\circ\)，所以 \(0^\circ<A<60^\circ\)，\(A\) 必为锐角. 结合题给
\(\sin21^\circ47'=\frac{3\sqrt3}{14}\)，得 \(A=21^\circ47'\). 我舰艇从 \(A\) 驶向 \(B\) 的方位角等于从 \(A\) 指向 \(C\) 的方位角加上 \(\angle A\)，故所求方位角为
\(45^\circ+21^\circ47'=66^\circ47'\).
接着 \(B=180^\circ-120^\circ-21^\circ47'=38^\circ13'\). 再由正弦定理，
\(\frac{9x}{\sin21^\circ47'}=\frac{10}{\sin38^\circ13'}\)，所以
\(x=\frac{10\sin21^\circ47'}{9\sin38^\circ13'}=\frac{10\cdot(3\sqrt3/14)}{9\cdot(5\sqrt3/14)}=\frac23\text{ 小时}=40\text{ 分钟}\).
因此我舰艇应沿 \(66^\circ47'\) 的方位角前去追捕，需用 \(\frac23\) 小时在 \(B\) 处追上敌舰.

【余弦定理核验】同一追捕时间还必须满足余弦定理
\((21x)^2=10^2+(9x)^2-2\cdot10\cdot9x\cos120^\circ\). 代入 \(\cos120^\circ=-\frac12\)，化简得
\(441x^2=100+81x^2+90x\)，即 \(36x^2-9x-10=0\). 因式分解为
\((3x-2)(12x+5)=0\)，故 \(x=\frac23\) 或 \(x=-\frac5{12}\). 由于 \(x>0\)，只取 \(x=\frac23\)，与上面的正弦定理结果一致.

【法二】设我舰艇在某处 \(B\) 追上敌舰需用 \(x\) 小时，则在 \(\triangle ABC\) 中，
\( AB=21x\text{ 海里} \)，\( BC=9x\text{ 海里} \).
又 \( AC=10\text{ 海里} \)，
\( \angle C=45^\circ+\left( 180^\circ-105^\circ \right)=120^\circ \).
根据正弦定理，得
\( \frac{21x}{\sin 120^\circ}=\frac{9x}{\sin A} \).
所以
\(\sin A=\frac{9x}{21x}\cdot \sin 120^\circ=\frac{3}{7}\cdot \sin 60^\circ=\frac{3\sqrt{3}}{14}\).
因为 \( \sin 21^\circ47'=\frac{3\sqrt{3}}{14} \)，又 \(A\) 为锐角，所以 \( \angle A=21^\circ47' \).
所以所求的方位角是 \( 45^\circ+21^\circ47'=66^\circ47' \).

根据余弦定理，得
\( \left( 21x \right)^2=10^2+\left( 9x \right)^2-2\times 10\times \left( 9x \right)\times \cos 120^\circ \).
整理，得
\( 36x^2-9x-10=0 \).
解这个方程，得
\( x_1=\frac{2}{3} \)，\( x_2=-\frac{5}{12} \)（不合题意，舍去）.

答：我舰艇应沿 \(66^\circ47'\) 的方位角前去追捕，需用 \(\frac23\) 小时，即 40 分钟.

【法三】
设我舰艇在某处 \(B\) 追上敌舰需用 \( x \) 小时，则在 \( \triangle ABC \) 中，
\( AB=21x\text{ 海里} \)，\( BC=9x\text{ 海里} \).
又 \( AC=10\text{ 海里} \)，
\( \angle C=45^\circ+\left( 180^\circ-105^\circ \right)=120^\circ \).
根据余弦定理，得
\( \left( 21x \right)^2=10^2+\left( 9x \right)^2-2\times 10\times \left( 9x \right)\times \cos 120^\circ \).
整理，得
\( 36x^2-9x-10=0 \).
解这个方程，得 \( x_1=\frac{2}{3} \)，\( x_2=-\frac{5}{12} \)（不合题意，舍去）.
因此
\( AB=21\times \frac{2}{3}=14\text{ 海里} \)，
\( BC=9\times \frac{2}{3}=6\text{ 海里} \).
又根据余弦定理，得
\( \cos A=\frac{10^2+14^2-6^2}{2\times 10\times 14}=\frac{13}{14} \).
因为 \( \sin 21^\circ47'=\frac{3\sqrt{3}}{14} \)，所以
\( \cos 21^\circ47'=\sqrt{1-\left( \frac{3\sqrt{3}}{14} \right)^2}=\frac{13}{14} \).
所以 \( \angle A=21^\circ47' \).
所以所求的方位角是 \( 45^\circ+21^\circ47'=66^\circ47' \).

答：我舰艇应沿 \(66^\circ47'\) 的方位角前去追捕，需用 \(\frac23\) 小时，即 40 分钟.
\end{solution}

\section{参考题}

\begin{problem}
某生产队要建造一个已知体积为 \(V\) 的有盖圆柱形氨水池，问这个氨水池的高 \( h \) 和底面半径 \( r \) 取多大时，用料最省？
\end{problem}

\begin{solution}
设 \(V>0\)，氨水池底面半径 \(r>0\)，高 \(h>0\). 有盖圆柱的用料面积就是表面积
\(S=2\pi r^2+2\pi rh\). 体积条件为
\(V=\pi r^2h\)，所以 \(h=\frac{V}{\pi r^2}\).

\solutionfigure{\bitmapfigure[max width=4.6cm]{img/1977/liaoning/q12_solution_fig1.png}}

代入表面积式，得
\(S=S(r)=2\pi r^2+\frac{2V}{r}\)，其中 \(r>0\). 求导得
\(S'(r)=4\pi r-\frac{2V}{r^2}=\frac{4\pi r^3-2V}{r^2}\). 因为分母 \(r^2>0\)，而 \(4\pi r^3-2V\) 随 \(r\) 严格递增，并在
\(r_0=\sqrt[3]{\frac{V}{2\pi}}\) 处为零，所以当 \(0<r<r_0\) 时 \(S'(r)<0\)，当 \(r>r_0\) 时 \(S'(r)>0\). 因而 \(S\) 先减后增，在且仅在 \(r=r_0\) 时取得最小值. 还可写为
\(r=\sqrt[3]{\frac{V}{2\pi}}=\frac{\sqrt[3]{4\pi^2V}}{2\pi}\).
此时
\(h=\frac{V}{\pi r^2}=\frac{2\pi r^3}{\pi r^2}=2r=\frac{\sqrt[3]{4\pi^2V}}{\pi}\).

答：当氨水池的高
\( h=\frac{\sqrt[3]{4\pi^2V}}{\pi} \)，
底面半径
\( r=\frac{\sqrt[3]{4\pi^2V}}{2\pi} \)，
即高和底面直径相等时，用料最省.
\end{solution}

\begin{problem}
求 \( y=\sin x \)（\( x \) 在 \( [0,\pi] \) 上）与 \( x \) 轴所围成的图形的面积.

\bitmapfigure[max width=6.0cm]{img/1977/liaoning/q13_fig1.png}
\end{problem}

\begin{solution}
因为 \(x\in[0,\pi]\) 时 \(\sin x\geqslant0\)，曲线 \(y=\sin x\) 位于 \(x\) 轴上方，所以所求面积等于定积分
\(S=\int_0^\pi\sin x\,\mathrm{d}x\). 由于 \((-\cos x)'=\sin x\)，有
\(S=\left[-\cos x\right]_0^{\pi}=-\cos\pi-(-\cos0)=1-(-1)=2\).
\end{solution}
